\section{Contexto y motivación: Relevancia de la detección precoz en enfermedades neurodegenerativas}

Este proyecto nace de la necesidad de abordar uno de los principales desafíos que enfrenta la medicina moderna: lograr un diagnóstico temprano, no invasivo, accesible y de bajo coste para patologías altamente prevalentes y con un profundo impacto en la calidad de vida de los pacientes. En particular, este trabajo se focaliza en las enfermedades neurodegenerativas, afecciones cuya incidencia se encuentra en continuo aumento debido al envejecimiento demográfico de la población mundial. Estas patologías se caracterizan por la pérdida progresiva de la función neuronal, conduciendo en la gran mayoría de los casos a un severo deterioro cognitivo, así como a una drástica disminución de la funcionalidad y de la independencia personal.

La detección precoz de estas alteraciones es crucial, ya que permite una intervención temprana y la implementación de estrategias de manejo clínico que pueden ralentizar la progresión de la enfermedad, mejorar el bienestar del paciente y brindar apoyo a su entorno familiar \cite{abrilcarreres2024}.

Sin embargo, los métodos de diagnóstico actuales presentan barreras de accesibilidad significativas. Técnicas como la resonancia magnética o el análisis de líquido cefalorraquídeo, si bien son clínicamente efectivas, resultan invasivas y requieren de infraestructuras complejas y personal altamente especializado. Esta dependencia de recursos clínicos avanzados genera limitaciones de disponibilidad geográfica, provocando que el acceso a un diagnóstico temprano sea desigual y dependa de la capacidad de la red sanitaria territorial \cite{avance2023}.

En este contexto, surge la necesidad de desarrollar herramientas de cribado alternativas. Este trabajo propone explorar el potencial de la Inteligencia Artificial y el análisis de datos para identificar patrones sutiles en la voz de los pacientes, tales como indicadores de disfonía, inestabilidad en la fonación o temblor neuromotor. Estas anomalías vocales podrían actuar como biomarcadores digitales indicativos de la presencia de enfermedades neurodegenerativas en sus fases iniciales. La fonación requiere una coordinación neuromuscular extremadamente precisa, convirtiendo a la voz en una fuente rica en información y en un candidato altamente prometedor para el diagnóstico temprano \cite{romero2018}.

Por consiguiente, la motivación central de este proyecto es aportar una solución tecnológica orientada a la democratización del prediagnóstico. El objetivo último es desarrollar y validar un sistema computacional que, analizando la voz como biomarcador, actúe como sistema de apoyo a la decisión clínica, ofreciendo a los pacientes una vía accesible para la identificación precoz de estas patologías.
